\printedpage{39}
% head: A. Grothendieck : Dérivateurs, Ch. II, Édition des Universités de Montpellier II et Paris VII

$I$ satisfying condition a$''$ (self-dual), condition b$'$ (also self-dual) is also satisfied (by case $4^o$ above). I say therefore that there exist $I$ and $i$ in $I$, such that
\[
\beta_{i,*} : (i\backslash I)^\wedge \longrightarrow I^\wedge
\]
is not conservative. (\textbf{NB} On the other hand $\alpha_{i,!} : (I/i)^\wedge \longrightarrow I^\wedge$ is always conservative, which gives back case $4^o$ as a special case of $3^o$). As $\beta_i : i\backslash I \longrightarrow I$ is $W_0$-proper (and even $W_\omega$-proper!), $\beta_{i,*}$ ``is computed fiber by fiber'', so we have
\msp{page 64}
\[
\beta_{i,*}(G)(j) \simeq \prod_{\alpha\in\mathrm{Hom}(i,j)} G(\alpha),
\]
the fiber in question at $j$ being discrete. Let us take
\[
I = 0 \underset{\beta}{\overset{\alpha}{\rightrightarrows}} 1, \qquad\quad \text{with } \alpha \neq \beta,
\]
then $0\backslash I$ is $\Psi = {}$\begin{tikzcd}[ampersand replacement=\&,row sep=0.1em, column sep=1.2em, baseline=(\tikzcdmatrixname-2-1.base)]
 \& 1 \\
0 \ar[ur, "\alpha'"] \ar[dr, "\beta'"'] \& \\
 \& 1'
\end{tikzcd}, and $\Psi^\wedge$ is the category of diagrams of sets $G = {}$\begin{tikzcd}[ampersand replacement=\&,row sep=0.1em, column sep=1.2em, baseline=(\tikzcdmatrixname-2-1.base)]
 \& G_1 \ar[dl, "\lambda"'] \\
G_0 \& \\
 \& G_{1'} \ar[ul, "\mu"]
\end{tikzcd}. The functor $\beta_*$ deduced from $\beta : i\backslash I \longrightarrow I$ transforms this diagram into
\[
G_1 \times G_{1'} \underset{\lambda\,\mathrm{pr}_2}{\overset{\lambda\,\mathrm{pr}_1}{\rightrightarrows}} G_0.
\]
\texttt{[rather} $\mu\,\mathrm{pr}_2$\texttt{]}. If we restrict ourselves to the diagrams $G$ for which $G_{1'} = \emptyset$, $\beta_*$ on these gives
\[
\emptyset \rightrightarrows G_0,
\]
so the restriction of $\beta_*$ to this subcategory (equivalent to $(\boldsymbol{\Delta}^1)^\wedge$) is not faithful, nor conservative, taking account only of the component $G_0$, and not of $G_1$.

\msp{page 63bis}

\textbf{Definition 4.6.} \textit{We say that $u : I' \longrightarrow I$ in} Diag \textit{is} $\D$-biaspheric \textit{if it satisfies the equivalent conditions b, b$'$, b$''$ of 4.4,} $\D$-connected \textit{if it satisfies condition a, or again (subject to the existence of $u_!$ resp. of $u_*$) the equivalent conditions a$'$, a$''$.} % [sic] Definition -> Définition, accent missing in print

Thus, $\D$-biaspheric implies $\D$-connected, and the converse is true when $I$ satisfies (with respect to $\D$) condition $3^o$ of 4.4 (thus also, if it satisfies one of the weaker conditions $2^o$, $1^o$), or the condition dual to this one (since the notions $\D$-biaspheric, $\D$-connected are self-dual). In this regard, let us note the

\textbf{Corollary 4.7.} \textit{Suppose that $\mathcal{A} = \D(e)$ satisfies the following condition: the functor}
\[
\left\{
\begin{array}{lcccl}
\mathcal{A} & \times & \mathcal{A} & \longrightarrow & \mathcal{A} \\[1ex]
(\xi & , & \eta) & \longrightarrow & \xi \amalg \eta
\end{array}
\right.
\]
\textit{is conservative (resp. faithful). Then for every étale arrow $u : J \longrightarrow I$ in} Diag \textit{(and in particular for $I/i \longrightarrow I$), the functor $u_! : \D(J) \longrightarrow \D(I)$ is conservative (resp. faithful)} (\footnote{\textbf{NB} This condition is satisfied if $\forall\, \xi, \eta \in \mathcal{A}$ we have $\mathrm{Hom}(\xi,\eta) \neq \emptyset$, (for it reduces to this: if $u_1$,
% footnote continues at the foot of p. 40
$u_2$ are two arrows in Ens between sets \textit{that are nonempty,} and if $u_1 \times u_2$ is \texttt{[an]} isomorphism, $u_1$, $u_2$ are \texttt{[some]} isomorphisms). This applies in particular if $\mathcal{A}$ admits a zero object, e.g. if $\mathcal{A}$ is additive.}). \textit{If this condition, or the dual condition,}
\printedpage{40}
% head: A. Grothendieck : Dérivateurs, Ch. II, Édition des Universités de Montpellier II et Paris VII
\msp{page 64bis}
\textit{involving the products $\xi\times\eta$ ($\xi$, $\eta$ in $\mathcal{A}$), is satisfied, then for every arrow $u : I' \longrightarrow I$ in} Diag, \textit{$u$ is $\D$-connected if and only if $u$ is $\D$-biaspheric.}

\textsc{Proof.} The hypothesis of conservativity, or of faithfulness carries over to sums indexed by an arbitrary set (subject to existence), as one sees by writing $\coprod\limits_{\lambda} \xi_\lambda = \xi_{\lambda_0} \amalg \coprod\limits_{\lambda\neq\lambda_0} \xi_\lambda$. The assertion made about $u_!$ then follows from the computation of $u_!$ fiber by fiber (which is valid since $u$ is étale hence smooth), taking into account that the fibers are discrete.

I know of no example of a pair $(\D, u)$ of a derivator $\D$ and an arrow $u$ in Diag, for which $u$ is $\D$-connected without being $\D$-aspheric. We have given conditions on $\D$, which ensure equivalence of the two notions. If $\D$ is of the form $\D_{\mathcal{M}}$, and if $\mathcal{M}$ satisfies the condition stated in 4.7 or its dual, then by 4.7 it is OK. So a counterexample with a $\D_{\mathcal{M}}$ would have to take this into account. I am not going to persist in trying to clear up how things stand.

\msp{page 65}

\textbf{Corollary 4.8.} \textit{An arrow $u$ of} Diag \textit{which is a universal $\D$-equivalence is $\D$-biaspheric, and \textup{a fortiori} $\D$-connected. Conversely, if $u$ is either $\D$-smooth, or $\D$-proper, or a weak $\D$-fibration, and if $u$ is $\D$-biaspheric, it is a universal $\D$-equivalence.}

The proof is that of the special case 4.3. We thus have a diagram of implications:
\[
\begin{tikzcd}[column sep=small, row sep=large]
\text{Cat-equivalence} \ar[dr, Rightarrow] & \D\text{-universal equivalence } (W_{\mathrm{univ}}) \ar[d, Rightarrow, "{\text{\footnotesize\begin{tabular}{@{}l@{}}strict implication ; equivalence\\ in the case where $u$ is either $\mathbf{D}$-proper,\\ or $\mathbf{D}$-smooth, or a weak $\mathbf{D}$-fibration.\end{tabular}}}"] & \\
 & \D\text{-biaspheric} \ar[dl, Rightarrow] \ar[d, Rightarrow, "{\text{\footnotesize\begin{tabular}{@{}l@{}}equivalence\\ in all\\ known cases\end{tabular}}}"] \ar[dr, Rightarrow] & \\
\D\text{-aspheric} \ar[dr, Rightarrow] & \D\text{-connected} \ar[d, Rightarrow] & \D\text{-coaspheric} \ar[dl, Rightarrow] \\
 & \D\text{-equivalences } (W) &
\end{tikzcd}
\]
(\footnote{\textbf{NB} These seven classes of morphisms in Cat are stable under composition except for the notion of $\D$-connectedness, all these notions depend only on the set $W = W_{\D}$ of $\D$-equivalences.}). Recall that for a given $u : I' \longrightarrow I$
\begin{enumerate}[leftmargin=*, labelindent=1.5em]
\item[a)] $u$ is $\D$-aspheric if and only if $\forall\, F$ \textit{constant} in $\D(I)$
\[
F \xrightarrow[\sim]{\text{adj.}} u_*u^*F \qquad\quad \text{is \texttt{[an]} isomorphism.} \tag*{$(\ast)$}
\]
\printedpage{41}
% head: A. Grothendieck : Dérivateurs, Ch. II, Édition des Universités de Montpellier II et Paris VII
\item[b)] $u$ is $\D$-coaspheric if and only if $\forall\, F$ \textit{constant} in $\D(I)$
\[
F \xleftarrow[\sim]{\text{adj.}} u_!u^*F \qquad\quad \text{is \texttt{[an]} isomorphism.} \tag*{$(\ast\ast)$}
\]
\item[c)] $u$ is $\D$-connected if and only if for \textit{every} $F$ in $\D(I)$ (not necessarily constant, nor locally constant), $(\ast)$ is \texttt{[an]} isomorphism, or what amounts to the same thing, for every $F$ $(\ast\ast)$ is \texttt{[an]} isomorphism.
\end{enumerate}

\msp{page 66}

Recall that the class $W_{\mathrm{asph}}$ \texttt{[of]} $W$-aspheric arrows and that, $\mathrm{Liss}_W$, of $W$-smooth arrows, determine each other, thus
\[
\left\{
\begin{array}{lcl}
(f : X \longrightarrow Y) \in W_{\mathrm{asph}} & \Longleftrightarrow & \forall\, (Y' \xrightarrow{\ u\ } Y) \in \mathrm{Liss}_W, \\
 & & X' = X \times_Y Y' \longrightarrow X' \text{ is } \in W. % [sic] X' = X \times_Y Y' \longrightarrow X' -> X' = X \times_Y Y' \longrightarrow Y', the target must be Y' (compare (4.12))
\\
(f : Y' \longrightarrow Y) \in \mathrm{Liss}_W & \Longleftrightarrow & \text{The functor } X \longmapsto X' = X \times_Y Y' \\
 & & \text{from } \mathrm{Cat}/Y \text{ to } \mathrm{Cat}/Y' \text{ transforms} \\
 & & \text{arrows in } W_{\mathrm{asph}} \text{ in \texttt{[rather into]}} \\
 & & \text{arrows in } W_{\mathrm{asph}}.
\end{array}
\right.
\tag*{(4.10)}
\]

To every subset $A$ of Fl(Cat) stable under isomorphisms, one can associate a subset $\mathbf{L}(A)$ of Fl(Cat), by the relation
\[
\begin{array}{rcl}
(u : Y' \longrightarrow Y) \in \mathbf{L}(A) & \Longleftrightarrow & u^* : \mathrm{Cat}/Y \longrightarrow \mathrm{Cat}/Y' \\
 & & \text{transforms arrows} \\
 & & \in A \text{ into arrows } \in A.
\end{array}
\tag*{(4.11)}
\]

We see that $\mathbf{L}(A)$ contains the isomorphisms, is stable under composition, and is also stable under base change (as is the case for $\mathrm{Liss}_W$). Conversely, to a subset $L$ of Fl(Cat), stable under isomorphisms, we associate a subset $\mathbf{A}_W(L)$ by
\[
\begin{array}{rcl}
(f : X \longrightarrow Y) \in \mathbf{A}_W(L) & \Longleftrightarrow & \forall\, (Y' \longrightarrow Y) \text{ in } L, \\
 & & X' = X \times_Y Y' \longrightarrow Y' \\
 & & \text{is in } W.
\end{array}
\tag*{(4.12)}
\]

If $L$ is stable under base change, then $\mathbf{A}_W(L)$ is stable under composition. It contains the isomorphisms (since $W$ contains them), is contained in $W$ if $L$ contains the identities. But
\msp{page 67}
no reason for it to be stable under arbitrary base change (and in fact, $W_{\mathrm{asph}}$ is not).

(\footnote{Botched \texttt{[up to the end of} \S4\texttt{]}, cf. \S6.}). We have, for a given $A \subseteq W \subseteq \mathrm{Fl(Cat)}$, stable under isomorphisms,
\[
A \subseteq \mathbf{A}_W(\mathbf{L}(A)) \overset{\text{def}}{=} A^\natural \subseteq W, \qquad \text{\footnotesize\begin{tabular}{@{}l@{}}subset stable under composition,\\ contains \texttt{[the]} isomorphisms,\end{tabular}}
\tag*{(4.13)}
\]
\printedpage{42}
% head: A. Grothendieck : Dérivateurs, Ch. II, Édition des Universités de Montpellier II et Paris VII
and for every subset $L$ of Fl(Cat) which is stable under base change, and under composition
\[
L \subseteq \mathbf{L}(\mathbf{A}_W(L)) \overset{\text{def}}{=} L^\flat, \qquad \text{\footnotesize\begin{tabular}{@{}l@{}}subset of Fl\,Cat stable\\ under composition and under base\\ change, containing the isomorphisms.\end{tabular}}
\tag*{(4.14)}
\]
We have
\[
\left\{
\begin{array}{lcl}
A^{\natural\natural} & = & A^\natural, \quad \text{since } \mathbf{A}_W(L)^\natural = \mathbf{A}_W(L) \\[1ex]
L^{\flat\flat} & = & L^\flat, \quad \text{since } \mathbf{L}(A)^\flat = \mathbf{L}(A).
\end{array}
\right.
\]
Examples:
\begin{enumerate}[leftmargin=*, labelindent=1em]
\item[$1^o$)] $L = \text{Ét}$, $\mathbf{A}_W(L) = W_{\mathrm{asph}}$, $\mathbf{L}(\mathbf{A}_W(L)) = L^\flat = \mathrm{Liss}_W$.
\item[$2^o$)] $L = \text{Coét}$, $\mathbf{A}_W(L) = W_{\mathrm{coas}}$, $\mathbf{L}(\mathbf{A}_W(L)) = L^\flat = \mathrm{Prop}_W$.
\item[$3^o$)] $A = W$, $\mathbf{L}(A) = \mathrm{Fib}_W$, $A^\natural = W$.
\item[$4^o$)] $L = \mathrm{Fl\,Cat}$, $\mathbf{A}_W(L) = W_{\mathrm{univ}}$, $L^\flat = L$.
\end{enumerate}
\msp{page 68 vide}

\clearpage
\printedpage{43}
% head: A. Grothendieck : Dérivateurs, Ch. II, Édition des Universités de Montpellier II et Paris VII
\msp{page 69}

\section*{5\quad Formal duality}

Let $\mathcal{P}$, $\mathcal{Q}$ be two ordered sets. I am trying to isolate a notion of ``duality'' (formal) linking $\mathcal{P}$ and $\mathcal{Q}$. This will be a relation between $\mathcal{P}$ and $\mathcal{Q}$
\[
\mathcal{R} \subseteq \mathcal{P} \times \mathcal{Q}
\tag*{(5.1)}
\]
having the following properties
\begin{enumerate}[leftmargin=*, labelindent=0.5em]
\item[(D)] For every subset $A$ of $\mathcal{P}$, the set
\end{enumerate}
\[
A^\cup \overset{\text{def}}{=} \{ q \in \mathcal{Q} \mid (p,q) \in \mathcal{R} \text{ for every } p \in A,\ \text{ i.e. } A \times \{q\} \subseteq \mathbf{R} \} % [sic] \mathbf{R} -> \mathcal{R}, the relation is written \mathcal{R} everywhere else
\]
is an open subset of $\mathcal{Q}$ having a greatest element $\beta(A) : A^\cup = \mathcal{Q}/\beta(A)$. Dually, for every subset $B$ of $\mathcal{Q}$, the set
\[
B^\cap \overset{\text{def}}{=} \{ p \in \mathcal{P} \mid (p,q) \in \mathcal{R} \text{ for every } q \in B,\ \text{ i.e. } \{p\} \times B \subseteq \mathcal{R} \}
\]
is an open subset of $\mathcal{P}$ having a greatest element $\alpha(B)$, i.e. $B^\cap = \mathcal{P}/\alpha(B)$.

Most often, $\mathcal{P}$ and $\mathcal{Q}$ will be stable under arbitrary suprema (hence also under arbitrary infima). In this case, (D) is equivalent, on the side of the $A^\cup$ (for $A \in \mathfrak{P}(\mathcal{P})$) to this:
\begin{enumerate}[leftmargin=*, labelindent=0em]
\item[(D$_0$)] For every $p \in \mathcal{P}$, $\mathcal{R}(p) = \{ q \in Q \mid (p,q) \in \mathcal{R} \}$ is an open \textit{filtering} subset of $\mathcal{Q}$, or what amounts to the same thing, a nonempty subset ($\Longleftrightarrow$ containing the least element $\emptyset_{\mathcal{Q}}$ of $\mathcal{Q}$) stable under suprema of two elements. Moreover, $\mathcal{R}(p)$ is stable under filtering suprema, or what amounts to the same thing (being itself filtering), it contains a greatest element. It amounts % [sic] Q -> \mathcal{Q}, plain italic Q printed in "q \in Q" instead of the script \mathcal{Q}
\msp{page 70}
to the same thing to say that $\mathcal{R}(p)$ is an \textit{open subset} of $\mathcal{Q}$ stable under arbitrary suprema (which implies that it contains $\emptyset_{\mathcal{Q}}$). Dually, for every $q \in \mathcal{Q}$, ${}^{s}\mathcal{R}(q) = \{ p \in \mathcal{P} \mid (p,q) \in \mathcal{R} \}$ is an \textit{open} subset of $\mathcal{P}$, stable under arbitrary suprema, i.e. an open subset having a greatest element. % doubt resolved against PDF leaf 15: raised s, transpose relation
\end{enumerate}

In other words, when one postulates condition (D) for $A$ resp. $B$ reduced to one element, then it remains satisfied for every subset. This comes from the formulas
\[
\left\{
\begin{array}{lcl}
(\bigcup_i A_i)^\cup & = & \bigcap_i A_i^\cup \\[1ex]
(\bigcup_j B_j)^\cap & = & \bigcap_j B_j^\cap,
\end{array}
\right.
\tag*{(5.2)}
\]
and from the fact that the intersection of open subsets of the form $\mathcal{P}/\alpha_i$ in $\mathcal{P}$ is of the form $\mathcal{P}/\alpha$, where $\alpha = \inf_i \alpha_i$, and likewise in $\mathcal{Q}$. Thus we find
\[
\left\{
\begin{array}{lcl}
\beta(\bigcup_i A_i) & = & \inf_i \beta(A_i) \\[1ex]
\alpha(\bigcup_j B_j) & = & \inf_j \alpha(B_j).
\end{array}
\right.
\tag*{(5.3)}
\]